\documentclass{article}
\usepackage[utf8]{inputenc}
\usepackage[portuguese]{babel}
\usepackage{amsmath}
\usepackage{amssymb}
\usepackage{geometry}
\geometry{a4paper, margin=1in}

\title{Lista 2 - Sistemas Baseados em Conhecimento (MAC0444/MAC5778)}
\author{Leonardo Heidi Almeida Murakami \\ NUSP: 11260186}
\date{\today}

\begin{document}

\maketitle

\section*{1. Exercicio 1}

Primeiro, formalizamos as sentenças dadas usando os predicados fornecidos:
\begin{itemize}
\item "Quem fez os exercícios vai bem na prova."  \\
    $\forall x (\text{fezEx}(x) \rightarrow \text{vaiBem}(x))$
\item "Quem vai bem na prova fica com média alta."  \\
    $\forall x (\text{vaiBem}(x) \rightarrow \text{mediaAlta}(x))$
\item "Quem fica com média alta é aprovado em mac444."  \\
    $\forall x (\text{mediaAlta}(x) \rightarrow \text{aprovado}(x))$
\item "João fez os exercícios."  \\
    $\text{fezEx}(\text{joao})$
\item "Maria foi bem na prova."  \\
    $\text{vaiBem}(\text{maria})$
\end{itemize}

Para usar a resolução SLD, convertemos as implicações para a forma clausal (cláusulas de Horn):
\begin{enumerate}
    \item $\text{vaiBem}(x) \leftarrow \text{fezEx}(x)$
    \item $\text{mediaAlta}(x) \leftarrow \text{vaiBem}(x)$
    \item $\text{aprovado}(x) \leftarrow \text{mediaAlta}(x)$
    \item $\text{fezEx}(\text{joao}).$
    \item $\text{vaiBem}(\text{maria}).$
\end{enumerate}

\subsection*{Prova para João}
Queremos provar $\text{aprovado}(\text{joao})$. Partindo da negação:
\begin{itemize}
    \item[G0:] $\leftarrow \text{aprovado}(\text{joao})$
    \item[G1:] Resolvendo com (3) e $\{x/\text{joao}\}$: \\
    $\leftarrow \text{mediaAlta}(\text{joao})$
    \item[G2:] Resolvendo com (2) e $\{x/\text{joao}\}$: \\
    $\leftarrow \text{vaiBem}(\text{joao})$
    \item[G3:] Resolvendo com (1) e $\{x/\text{joao}\}$: \\
    $\leftarrow \text{fezEx}(\text{joao})$
    \item[G4:] Resolvendo com o fato (4), obtemos $\square$.
\end{itemize}
Como derivamos a cláusula vazia, temos que $\text{aprovado}(\text{joao})$ é consequência lógica da base.

\subsection*{Prova para Maria}
Queremos provar $\text{aprovado}(\text{maria})$. Partindo da negação:
\begin{itemize}
    \item[G0:] $\leftarrow \text{aprovado}(\text{maria})$
    \item[G1:] Resolvendo com (3) e $\{x/\text{maria}\}$: \\
    $\leftarrow \text{mediaAlta}(\text{maria})$
    \item[G2:] Resolvendo com (2) e $\{x/\text{maria}\}$: \\
    $\leftarrow \text{vaiBem}(\text{maria})$
    \item[G3:] Resolvendo com o fato (5), obtemos $\square$.
\end{itemize}
Como derivamos a cláusula vazia, temos que $\text{aprovado}(\text{maria})$ é consequência lógica da base.

\section*{2. Exercicio 2}

Vamos mostrar que o encadeamento para trás produz SIM para o objetivo $P(a)$ usando a base de conhecimento dada.

Base de conhecimento:
\begin{enumerate}
    \item $P(x) \leftarrow A1(x), A2(x)$
    \item $A1(x) \leftarrow B1(x), B2(x)$
    \item $A2(x) \leftarrow B3(x), B4(x)$
    \item $B1(a).$
    \item $B2(a).$
    \item $B3(a).$
    \item $B4(a).$
\end{enumerate}

Queremos provar $P(a)$:
\begin{itemize}
    \item[G0:] $\leftarrow P(a)$
    \item[G1:] Resolvendo com (1) e $\{x/a\}$: \\
    $\leftarrow A1(a), A2(a)$
    \item[G2:] Resolvendo $A1(a)$ com (2) e $\{x/a\}$: \\
    $\leftarrow B1(a), B2(a), A2(a)$
    \item[G3:] Resolvendo $B1(a)$ com o fato (4): \\
    $\leftarrow B2(a), A2(a)$
    \item[G4:] Resolvendo $B2(a)$ com o fato (5): \\
    $\leftarrow A2(a)$
    \item[G5:] Resolvendo $A2(a)$ com (3) e $\{x/a\}$: \\
    $\leftarrow B3(a), B4(a)$
    \item[G6:] Resolvendo $B3(a)$ com o fato (6): \\
    $\leftarrow B4(a)$
    \item[G7:] Resolvendo $B4(a)$ com o fato (7): \\
    $\leftarrow$ (lista vazia)
\end{itemize}
Como todos os sub-objetivos foram satisfeitos, $P(a)$ foi provado. Logo, a resposta é \textbf{SIM}.

\newpage
\section*{3. Exercicio 3}

\subsection*{(a) Resposta da consulta}
Programa Prolog:
\begin{verbatim}
result([_, E | L], [E | M]):- !, result(L, M).
result(_, []).
\end{verbatim}

Execução de \texttt{?- result([a, b, c, d, e, f, g], X).}:
\begin{enumerate}
    \item \texttt{result([a, b, c, d, e, f, g], X)} casa com a primeira regra:
    \begin{itemize}
        \item \texttt{\_} = \texttt{a}, \texttt{E} = \texttt{b}, \texttt{L} = \texttt{[c, d, e, f, g]}
        \item \texttt{X} = \texttt{[b | M]}
        \item Chamada recursiva: \texttt{result([c, d, e, f, g], M)}
    \end{itemize}
    \item \texttt{result([c, d, e, f, g], M)} casa com a primeira regra:
    \begin{itemize}
        \item \texttt{\_} = \texttt{c}, \texttt{E} = \texttt{d}, \texttt{L} = \texttt{[e, f, g]}
        \item \texttt{M} = \texttt{[d | M1]} (logo \texttt{X} = \texttt{[b, d | M1]})
        \item Chamada recursiva: \texttt{result([e, f, g], M1)}
    \end{itemize}
    \item \texttt{result([e, f, g], M1)} casa com a primeira regra:
    \begin{itemize}
        \item \texttt{\_} = \texttt{e}, \texttt{E} = \texttt{f}, \texttt{L} = \texttt{[g]}
        \item \texttt{M1} = \texttt{[f | M2]} (logo \texttt{X} = \texttt{[b, d, f | M2]})
        \item Chamada recursiva: \texttt{result([g], M2)}
    \end{itemize}
    \item \texttt{result([g], M2)} não casa com a primeira regra (lista com só um elemento).
    \item Casa com a segunda regra \texttt{result(\_, [])}:
    \begin{itemize}
        \item \texttt{M2} = \texttt{[]}
    \end{itemize}
    \item Voltando nas chamadas: \texttt{M2} = \texttt{[]}, então \texttt{M1} = \texttt{[f]}, \texttt{M} = \texttt{[d, f]} e \texttt{X} = \texttt{[b, d, f]}.
\end{enumerate}
Resposta:
\begin{verbatim}
X = [b, d, f].
\end{verbatim}

\subsection*{(b) Descrição do programa}

\textbf{O que o programa faz:}

O programa \texttt{result/2} retorna os elementos nas posições pares (2, 4, 6, ...) da lista de entrada. A cada chamada recursiva, ele pula o primeiro elemento e pega o segundo, continuando com o resto da lista.

\textbf{Casos cobertos pelo fato:}

O fato \texttt{result(\_, [])} é o caso base que cobre:
\begin{itemize}
    \item Lista vazia: \texttt{[]}.
    \item Lista com um único elemento: ex. \texttt{[g]}.
\end{itemize}
Em ambos os casos, a primeira regra não se aplica (não há pelo menos dois elementos), então o resultado é a lista vazia.

\textbf{Efeito do corte (\texttt{!}):}

O corte impede backtracking após a primeira regra ser aplicada. Sem ele, o Prolog tentaria usar a segunda regra (\texttt{result(\_, [])}) na mesma chamada, gerando uma solução alternativa incorreta (\texttt{X = []}). O corte garante que, se a lista tem pelo menos dois elementos, só a primeira regra será usada.

\textbf{Uso da variável anônima (\texttt{\_}):}

\begin{itemize}
    \item Em \texttt{[\_, E | L]}: indica o primeiro elemento, que queremos ignorar. Precisamos dele na estrutura do pattern matching para acessar \texttt{E} (segundo elemento), mas seu valor não importa.
    \item Em \texttt{result(\_, [])}: representa qualquer lista de entrada. Como sempre retornamos \texttt{[]}, não importa o conteúdo da entrada.
\end{itemize}

\end{document}